\documentclass[11pt]{article}
\usepackage[T1]{fontenc}
\usepackage[margin=1in]{geometry}
\usepackage{booktabs,longtable,microtype,hyperref}
\hypersetup{hidelinks}
\setlength{\parindent}{0pt}
\setlength{\parskip}{6pt}
\setlength{\emergencystretch}{3em}
\title{Folio: Manuscript Assessment}
\author{Jev-assisted review}
\date{2026-09-19}
\begin{document}
\maketitle
\section{Manuscript}
Uncertainty-aware document classification
\par File: \texttt{structured-synthetic.pdf}
\par Review ID: \texttt{f83219bd-01f8-48c9-a504-dc7123a695eb}
\par SHA-256: \texttt{30bfd350be4f8c13c927a36e4367c3d2}\allowbreak\texttt{c6644f6e66aa198fce2629893e97796a}
Pages: 3. Review units: 8. Profile: research.
\section{Overall assessment}
Score: 77.4/100. Recommendation: Promising ; minor revisions.
Review status: Complete. Mean model confidence: 50.0\%.
\section{Method}
Each passage is evaluated on five five-level Jev Score questions. Scores are mapped from [0,4] to [0,100]. Clarity contributes 16\%, rigor 28\%, evidence 24\%, completeness 12\%, and venue fit 20\%. Long-section passage scores are weighted by character count. Section weights are fixed by role and divided within a role by word count. Mean confidence describes concentration of model outputs, not the probability that the paper is correct. Findings and suggestions below are rubric-derived templates, not generated expert explanations. Source excerpts are context, not verified rationales.
\section{Section scores}
\begin{longtable}{p{.44\textwidth}rrrr}
\toprule Section & Pages & Score & Conf. & Weight \\ \midrule\endhead
Abstract & 1--1 & 82.9 & 55\% & 5.0\% \\
1 Introduction & 1--1 & 72.6 & 40\% & 10.0\% \\
2 Related Work & 1--1 & 75.1 & 43\% & 10.0\% \\
3 Methods & 2--2 & 85.9 & 62\% & 25.0\% \\
4 Results & 2--2 & 66.2 & 41\% & 25.0\% \\
5 Discussion and Limitations & 2--2 & 86.8 & 56\% & 12.0\% \\
6 Conclusion & 2--2 & 82.6 & 56\% & 8.0\% \\
References & 3--3 & 69.4 & 39\% & 5.0\% \\
\bottomrule\end{longtable}
\section{Target venue and website grounding}
\textbf{NeurIPS 2026}
\par Supplied website: \url{https://neurips.cc/Conferences/2026/ReviewerGuidelines}
\par Track/category: General
\par Retrieved website passages inform all criteria. Venue fit contributes 20\% to the overall score. These are Folio's application weights, not the venue's official scale or an estimate of acceptance probability. The JSON export contains all selected passages.
\subsection*{2026 Reviewer Guidelines}
\url{https://neurips.cc/Conferences/2026/ReviewerGuidelines}
\par Retrieved: 2026-09-19T13:41:45.828780+00:00
\par Source SHA-256: \texttt{d7a667d7d1268b3976635ff5b912406d}\allowbreak\texttt{fcee92baba40a671777fa6c04c11d334}
\begin{quote}\small General Reviewing Guidelines / Quality: Is the submission technically sound? Are claims well supported (e.g., by theoretical analysis or experimental results)? Are the methods used ...\end{quote}
\section{Detailed findings}
\subsection{Abstract}
\par\textbf{Clarity (88.0/100):} Precise definitions and coherent organization make the meaning unambiguous.
\par\textbf{Rigor (85.0/100):} The reasoning is well justified with explicit assumptions and minor gaps.
\par\textbf{Evidence (86.2/100):} Central claims have specific relevant support with small omissions.
\par\textbf{Completeness (88.2/100):} The section's purpose is fully addressed with the details a critical reader needs.
\par\textbf{Venue fit (68.8/100):} The contribution aligns with the venue's documented scope and addresses the applicable scientific review priorities.
\par Revision prompt: Make the contribution's connection to the cited venue scope and applicable review priorities explicit.
\begin{quote}\small We study whether explicit uncertainty reporting improves the interpretability of a small
document-classification system. In this synthetic demonstration, a logistic regression baseline is
compared with temperature-scaled predictions on a fixed, balanced collection of 600 generated
abstracts. We report held-out accuracy and calibration error over five random seeds. Calibration
improves while accuracy is unchanged within sampling variation. These illustrative results describe a
controlled toy setting; they do not establish performance on real sch\end{quote}
Passages evaluated: 1. Model: jev-1.13.0.
\subsection{1 Introduction}
\par\textbf{Clarity (86.5/100):} The argument is easy to follow with only minor ambiguity.
\par\textbf{Rigor (85.2/100):} The reasoning is well justified with explicit assumptions and minor gaps.
\par\textbf{Evidence (50.8/100):} Some relevant support is provided but a central claim remains insufficiently supported.
\par Revision prompt: Connect each central claim to a specific result, derivation, or relevant citation.
\par\textbf{Completeness (79.2/100):} The necessary details are present with only minor omissions.
\par\textbf{Venue fit (66.0/100):} The contribution aligns with the venue's documented scope and addresses the applicable scientific review priorities.
\par Revision prompt: Make the contribution's connection to the cited venue scope and applicable review priorities explicit.
\begin{quote}\small Automated document triage often exposes a class label without the uncertainty behind it. Readers may
mistake a confident prediction for a reliable one. Our research question is whether a simple
post-processing calibration step can improve the match between prediction confidence and observed
accuracy without changing the classifier. The contribution is a small, reproducible evaluation protocol
separating training, calibration, and testing. This original synthetic manuscript is a software
demonstration and is not a published scientific study.\end{quote}
Passages evaluated: 1. Model: jev-1.13.0.
\subsection{2 Related Work}
\par\textbf{Clarity (84.5/100):} The argument is easy to follow with only minor ambiguity.
\par\textbf{Rigor (80.8/100):} The reasoning is well justified with explicit assumptions and minor gaps.
\par\textbf{Evidence (71.5/100):} Central claims have specific relevant support with small omissions.
\par Revision prompt: Connect each central claim to a specific result, derivation, or relevant citation.
\par\textbf{Completeness (54.5/100):} The main information is present but an important detail is missing.
\par Revision prompt: Add the missing details a reader would need to assess this section independently.
\par\textbf{Venue fit (76.5/100):} The contribution aligns with the venue's documented scope and addresses the applicable scientific review priorities.
\begin{quote}\small Proper scoring rules formalize evaluation of probabilistic predictions [1]. Calibration methods adjust
confidence estimates while preserving a predictor's ranking, with temperature scaling providing a simple
example [2]. Our protocol uses these established ideas rather than claiming a new calibration algorithm.
A broader comparison with isotonic regression would help establish whether the chosen baseline is
appropriate; that comparison is left for future work.
Folio / illustrative manuscript                                      1\end{quote}
Passages evaluated: 1. Model: jev-1.13.0.
\subsection{3 Methods}
\par\textbf{Clarity (94.8/100):} Precise definitions and coherent organization make the meaning unambiguous.
\par\textbf{Rigor (89.8/100):} The reasoning is internally sound and its assumptions and inference boundaries are explicit.
\par\textbf{Evidence (84.5/100):} Central claims have specific relevant support with small omissions.
\par\textbf{Completeness (77.2/100):} The necessary details are present with only minor omissions.
\par\textbf{Venue fit (80.2/100):} The contribution aligns with the venue's documented scope and addresses the applicable scientific review priorities.
\begin{quote}\small We construct 600 synthetic abstracts balanced across three topic classes. A fixed stratified split
allocates 360 abstracts to training, 120 to calibration, and 120 to testing. A TF-IDF representation with
unigrams and bigrams feeds a multinomial logistic regression with L2 regularization and C=1.
Vocabulary construction uses training data only. Temperature is selected on the calibration split by
minimizing negative log likelihood. The held-out test split is never used for model or temperature
selection. We repeat the procedure for random seeds \end{quote}
Passages evaluated: 1. Model: jev-1.13.0.
\subsection{4 Results}
\par\textbf{Clarity (93.5/100):} Precise definitions and coherent organization make the meaning unambiguous.
\par\textbf{Rigor (85.0/100):} The reasoning is well justified with explicit assumptions and minor gaps.
\par\textbf{Evidence (58.2/100):} Some relevant support is provided but a central claim remains insufficiently supported.
\par Revision prompt: Connect each central claim to a specific result, derivation, or relevant citation.
\par\textbf{Completeness (57.5/100):} The main information is present but an important detail is missing.
\par Revision prompt: Add the missing details a reader would need to assess this section independently.
\par\textbf{Venue fit (32.8/100):} The connection to the documented venue scope is weak and major scientific priorities are unaddressed.
\par Revision prompt: Make the contribution's connection to the cited venue scope and applicable review priorities explicit.
\begin{quote}\small In this synthetic example, the baseline achieves 81.2\% accuracy with a standard deviation of 1.8
percentage points across seeds. Temperature scaling achieves the same accuracy, as expected
because it preserves the maximal-logit class. Mean expected calibration error decreases from 0.142 to
0.061; standard deviations are 0.024 and 0.018 respectively. These numbers are illustrative fixture
values, not outputs of a conducted experiment. No statistical significance or real-world improvement is
claimed. A seed-level result table and confidence inter\end{quote}
Passages evaluated: 1. Model: jev-1.13.0.
\subsection{5 Discussion and Limitations}
\par\textbf{Clarity (87.8/100):} Precise definitions and coherent organization make the meaning unambiguous.
\par\textbf{Rigor (93.8/100):} The reasoning is internally sound and its assumptions and inference boundaries are explicit.
\par\textbf{Evidence (78.8/100):} Central claims have specific relevant support with small omissions.
\par\textbf{Completeness (91.5/100):} The section's purpose is fully addressed with the details a critical reader needs.
\par\textbf{Venue fit (83.2/100):} The contribution aligns with the venue's documented scope and addresses the applicable scientific review priorities.
\begin{quote}\small The protocol illustrates why discrimination and calibration should be reported separately. The apparent
calibration benefit applies only to the constructed distribution and should not be generalized to natural
abstracts. Generated text can encode label shortcuts, a small test split makes bin-based calibration
estimates unstable, and the experiment does not test distribution shift. The absence of executed
experiments is a central limitation of this demonstration. Future work should use independently
annotated documents, report inter-annotator ag\end{quote}
Passages evaluated: 1. Model: jev-1.13.0.
\subsection{6 Conclusion}
\par\textbf{Clarity (83.0/100):} The argument is easy to follow with only minor ambiguity.
\par\textbf{Rigor (95.8/100):} The reasoning is internally sound and its assumptions and inference boundaries are explicit.
\par\textbf{Evidence (81.8/100):} Central claims have specific relevant support with small omissions.
\par\textbf{Completeness (75.5/100):} The necessary details are present with only minor omissions.
\par\textbf{Venue fit (69.0/100):} The contribution aligns with the venue's documented scope and addresses the applicable scientific review priorities.
\par Revision prompt: Make the contribution's connection to the cited venue scope and applicable review priorities explicit.
\begin{quote}\small We describe a transparent protocol for separating classifier accuracy from calibration quality. The
synthetic values make the interface easy to inspect but supply no empirical evidence for deployment. A
substantive contribution would require execution of the protocol on real documents, external replication,
and careful analysis of uncertainty. Explicit scope and limitations are essential when presenting
model-based assessments.
Folio / illustrative manuscript                                      2\end{quote}
Passages evaluated: 1. Model: jev-1.13.0.
\subsection{References}
\par\textbf{Clarity (93.5/100):} Precise definitions and coherent organization make the meaning unambiguous.
\par\textbf{Rigor (32.5/100):} Central assumptions are unexamined and major inferential gaps remain.
\par Revision prompt: State the assumptions and justify the inferential steps that connect evidence to claims.
\par\textbf{Evidence (82.0/100):} Central claims have specific relevant support with small omissions.
\par\textbf{Completeness (87.5/100):} The section's purpose is fully addressed with the details a critical reader needs.
\par\textbf{Venue fit (75.8/100):} The contribution aligns with the venue's documented scope and addresses the applicable scientific review priorities.
\begin{quote}\small [1] Gneiting, T. and Raftery, A. E. (2007). Strictly proper scoring rules, prediction, and estimation. Journal
of the American Statistical Association, 102(477), 359-378.
[2] Guo, C., Pleiss, G., Sun, Y., and Weinberger, K. Q. (2017). On Calibration of Modern Neural
Networks. Proceedings of ICML, PMLR 70, 1321-1330.
Folio / illustrative manuscript                                      3\end{quote}
Passages evaluated: 1. Model: jev-1.13.0.
\section{Limitations and provenance}
\begin{itemize}
\item Grounded in retrieved guidance for NeurIPS 2026, track: General. Venue fit contributes 20\%; the other rubric dimensions contribute 80\%. These application weights are not the venue's official rating scale or acceptance probability.
\item Grounding uses selected passages from the supplied site, not an exhaustive policy or submission-compliance check.
\item This is a rubric-based assessment of extracted text, not a peer-review verdict. Figures, mathematical correctness, citation validity, and novelty require expert verification.
\item Recommendation thresholds and weights are design choices and have not been calibrated against expert review outcomes.
\item Non-ASCII characters are transliterated or replaced in this portable LaTeX export; use JSON for exact text.
\end{itemize}
Rubric: folio-1.1. Time: 1.28 seconds. Input tokens: 21518; output tokens: 600.
\par API documentation: \url{https://docs.typesafe.ai/primitives/score}
\end{document}
